\documentclass[11pt]{article}
\usepackage[T1]{fontenc}
\usepackage{lmodern,amsmath,amssymb,amsthm,mathtools}
\usepackage[margin=28mm]{geometry}
\usepackage[hidelinks]{hyperref}
\usepackage{microtype}
\newtheorem{theorem}{Theorem}[section]
\newtheorem{lemma}[theorem]{Lemma}
\newtheorem{proposition}[theorem]{Proposition}
\newtheorem{remark}[theorem]{Remark}
\newcommand{\R}{\mathbb R}
\newcommand{\C}{\mathbb C}
\newcommand{\T}{\mathbb T}
\newcommand{\rank}{\operatorname{rank}}
\newcommand{\ip}[2]{\langle #1,#2\rangle}
\title{Coarse information and finite rank obstructions\newline for flat lattice Hodge forms}
\author{Lluis Eriksson}
\date{30 October 2026}
\begin{document}
\maketitle
\begin{abstract}
On a finite periodic lattice with unit spacing, a flat Hodge form cannot have
a positive coercivity constant uniform under refinement if it is supplemented
only by observations in a fixed finite-dimensional space and uniformly bounded
rank fine-space repairs. The observations may vary arbitrarily with the
refinement, and their penalty need only vanish at zero. We prove the same
obstruction directly on kernels of uniformly bounded rank constraints.
The proof constructs a real physical Fourier subspace of dimension
$2(2R+1)^d$ with Hodge energy at most $4\pi^2dR^2/n^2$ times its squared norm.
Its intersection with the unobserved directions gives a necessary rank budget.
We also explain why this budget is not sufficient, using a two-dimensional
cross-coupling example and a classical complete criterion for a fixed map.
The Fourier construction, finite budget, and uniform obstructions have Lean
proofs; the fixed-map criterion is an analytic result here. These are finite
flat statements, with no interacting or continuum mass-gap claim.
\end{abstract}

\section{Purpose and relation to the existing record}
This text consolidates a continuation of the exact-kernel obstruction in
\cite{eriksson}. That paper exhibits fields in the kernel of a particular
line-integral block map. Changing the map can remove a particular witness.
Here the obstruction depends instead on the amount of coarse information:
a whole low-energy subspace is too large to be seen by a fixed-dimensional
measurement together with a bounded-rank repair. The map need not factor
through the original block variable, be local, or obey a norm bound.
The restricted-sector result is proved separately and does not follow merely
from failure of an estimate on the full space.

Fourier diagonalization, rank-nullity and finite-dimensional Schur complements
are classical. Finite-difference Hodge theory has a long history; Dodziuk's
work \cite{dodziuk} provides historical context, rather than this paper's
specific lattice normalization. We make no claim to a new general spectral
principle or an exhaustive priority search. The contribution is a concrete,
reproducibly checked constraint on proposed repair routes in this programme.
This consolidation is intended for the existing paper dossier, rather than
an additional deposit for each lemma.

\section{Spaces and the actual flat form}
Fix integers $d\geq1$, $n\geq1$ and $N_c\geq2$. Set
$X_n=(\mathbb Z/n\mathbb Z)^d$ and $m=N_c^2-1$. The real one-cochain space is
\[
 E_n=\{A:X_n\times\{1,\ldots,d\}\longrightarrow\R^m\},\qquad
 \|A\|^2=\sum_{x\in X_n}\sum_{i=1}^d\|A_i(x)\|^2.
\]
The internal coordinates are those of the repository's Euclidean Lie-coordinate
model. No identification with a matrix Lie algebra is required here.
For positive unordered plaquettes $i<k$, define
\begin{align*}
 (\mathrm dA)_{ik}(x)&=A_k(x+e_i)-A_k(x)-A_i(x+e_k)+A_i(x),\\
 (\delta A)(x)&=\sum_i\bigl(A_i(x)-A_i(x-e_i)\bigr).
\end{align*}
The flat operator $K_{0,n}$ is the actual repository Hodge operator, with
\begin{equation}\label{eq:hodge}
 \ip{A}{K_{0,n}A}=\|\mathrm dA\|^2+\|\delta A\|^2.
\end{equation}
The plaquette norm sums each pair once. All differences have unit spacing;
there is no hidden $n^2$ or block-side squared factor.

Choose a direction $j$ and two distinct internal unit coordinates $u,v$.
They exist since $m\geq3$. For $f:X_n\to\C$, put
\[
 (E_jf)_i(x)=\begin{cases}(\Re f(x))u+(\Im f(x))v,&i=j,\\0,&i\ne j.\end{cases}
\]
This is a real-linear injection and $\|E_jf\|^2=\sum_x|f(x)|^2$.
Writing $D_i f(x)=f(x+e_i)-f(x)$, the actual curl and divergence give
\begin{equation}\label{eq:stencil}
 \ip{E_jf}{K_{0,n}E_jf}
 =\sum_{i\ne j}\sum_x|D_i f(x)|^2+
   \sum_x|f(x)-f(x-e_j)|^2
 =\sum_i\sum_x|D_i f(x)|^2.
\end{equation}
The last equality is finite translation invariance. The exclusive two
plaquette incidences and the divergence term account for every direction
exactly once, with no extra factor two.

\section{A constructed low band}
For $R\geq0$ with $2R<n$, let $\Lambda_R=\{-R,\ldots,R\}^d$ and
$B=|\Lambda_R|=(2R+1)^d$. Define
\[
 \chi_k(x)=\exp(2\pi\mathrm i\,k\cdot x/n),\qquad
 f_a=\sum_{k\in\Lambda_R}a_k\chi_k,\qquad
 S_R=\{E_j f_a:a\in\C^{\Lambda_R}\}\subset E_n.
\]
\begin{lemma}\label{lem:band}
The real subspace $S_R$ has dimension $2B$, and for every $A\in S_R$,
\begin{equation}\label{eq:band}
 \ip{A}{K_{0,n}A}\leq b_R(n)\|A\|^2,
 \qquad b_R(n)=\frac{4\pi^2dR^2}{n^2}.
\end{equation}
\end{lemma}
\begin{proof}
For $k,\ell\in\Lambda_R$, the product root-of-unity sum is
$\sum_x\overline{\chi_k(x)}\chi_\ell(x)=n^d\mathbf1_{k=\ell}$.
The no-alias condition ensures that coordinate differences, of absolute value
at most $2R<n$, cannot be nonzero multiples of $n$. Thus synthesis is injective
and its complex image has dimension $B$. As a real space that image has
dimension $2B$; applying the real-linear injection $E_j$ preserves dimension.
The factor two represents independent real and imaginary coefficients stored
in two real internal coordinates, not two copies of a real scalar band.

For completeness, energy orthogonality is needed as well as mass orthogonality.
For each $i$,
\[
 D_i f_a=\sum_k a_k\bigl(e^{2\pi\mathrm i k_i/n}-1\bigr)\chi_k.
\]
Applying the same Gram sum to these derivative amplitudes gives
\begin{align*}
 \|E_jf_a\|^2&=n^d\sum_k|a_k|^2,\\
 \ip{E_jf_a}{K_{0,n}E_jf_a}
 &=n^d\sum_k |a_k|^2\sum_i|e^{2\pi\mathrm i k_i/n}-1|^2.
\end{align*}
Since $|e^{\mathrm it}-1|\leq|t|$ and $|k_i|\leq R$, each sum of multipliers
is at most $4\pi^2dR^2/n^2$. This proves the bound for every member of the
constructed image, not only its individual generators.
\end{proof}

\section{The finite information budget}
Let $F$ be a finite-dimensional real vector space, $C:E_n\to F$ a linear
measurement and $V:E_n\to E_n$ a linear operator. No symmetry or positivity
of $V$ is needed. Let $\Phi:F\to\R$ be any function with $\Phi(0)=0$.
\begin{theorem}\label{thm:budget}
If $c>b_R(n)$ and
\begin{equation}\label{eq:estimate}
 c\|A\|^2\leq\ip{A}{K_{0,n}A}+\Phi(CA)+\ip{A}{VA}
 \quad\text{for all }A\in E_n,
\end{equation}
then $2(2R+1)^d\leq\dim F+\rank V$.
\end{theorem}
\begin{proof}
If the budget is smaller than $\dim S_R$, restrict the map $(C,V)$ to $S_R$.
Its image has dimension at most $\dim F+\rank V$. Rank-nullity yields a
nonzero $A\in S_R$ with $CA=VA=0$. Equation~\eqref{eq:estimate} and
Lemma~\ref{lem:band} imply $c\|A\|^2\leq b_R(n)\|A\|^2$, a contradiction.
\end{proof}
The codomain dimension is a conservative measurement budget. We do not label
this formal endpoint as a sharper theorem with $\rank C$ in its place.
No regularity, sign, polynomial form or growth restriction on $\Phi$ enters.

\begin{proposition}\label{prop:sector}
If $P:E_n\to E_n$ is linear, $c>b_R(n)$ and
$\dim F+\rank P<2(2R+1)^d$, the inequality
\[
 c\|A\|^2\leq\ip{A}{K_{0,n}A}+\Phi(CA)
 \qquad(A\in\ker P)
\]
is false.
\end{proposition}
\begin{proof}
Apply the same rank-nullity argument to $(C,P)|_{S_R}$. Its nonzero kernel
vector belongs to the stipulated sector as well as the unobserved band.
\end{proof}
For example, $d=4,n=16,R=1$ gives dimension $162$ and
$b_R(n)=\pi^2/16<1$. Twelve measurement coordinates and a repair of rank at
most $149$ cannot yield~\eqref{eq:estimate} with $c=1$. This illustrates a
necessary budget, not the programme's chosen target constant.

\section{Uniform refinement and sector exclusion}
Fix $N'\geq1$ and put $n_L=(L+1)N'$ for $L\geq0$.
The finite-dimensional space $F$ is fixed. Allow arbitrary families
$C_L:E_{n_L}\to F$, $\Phi_L:F\to\R$ with $\Phi_L(0)=0$, and
$V_L:E_{n_L}\to E_{n_L}$ with $\rank V_L\leq r$ for one fixed integer $r$.
\begin{theorem}\label{thm:uniform}
There is no $c>0$ for which~\eqref{eq:estimate} holds for every $L$ and every
$A\in E_{n_L}$ with these families. Likewise, for linear $P_L$ of rank at
most $r$, there is no positive uniform $c$ for the sector estimates of
Proposition~\ref{prop:sector} on all $\ker P_L$.
\end{theorem}
\begin{proof}
Set $M=\dim F$ and $R=M+r+1$. Since $d\geq1$,
$2(2R+1)^d>M+r$. Choose $L$ so large that $2R<n_L$ and $b_R(n_L)<c$.
Here is a fully finite choice argument, without an assumed limit theorem:
with $a=4\pi^2dR^2$, take an integer
$L>\max\{2R,a/c,1\}$. Since $N'\geq1$, $n_L\geq L+1>L>1$, and
$a<cL\leq cn_L^2$. Theorem~\ref{thm:budget} now contradicts the rank bound.
Use Proposition~\ref{prop:sector} for the separate sector conclusion.
\end{proof}
The constants $c,r$ precede the quantifier over $L$. Arbitrary rescalings of
$C_L$ preserve its fixed codomain and do not escape the theorem. Removing
harmonic or constant directions is covered only when a concrete linear
constraint describing the proposed sector has uniformly bounded rank.
A quotient or its norm is not silently identified with such a sector.

\section{Why dimension is not a sufficiency criterion}
The following analytic statement is included to prevent the reverse inference
from a necessary rank budget. It is not a new Lean-verified endpoint here.
Let $E$ be a finite-dimensional real inner-product space, $H=H^*$ and
$B=H-cI$. Fix linear $C:E\to F$, and decompose $E=N\oplus W$ with
$N=\ker C$, $W=N^\perp$. In this decomposition write
\[
 B=\begin{pmatrix}A&D\\D^*&G\end{pmatrix}.
\]
\begin{proposition}\label{prop:schur}
A finite real-valued function $\Phi$ with $\Phi(0)=0$ can make
$\ip{x}{Bx}+\Phi(Cx)\geq0$ for every $x$ if and only if
$A\succeq0$ and $\ker A\subseteq\ker D^*$. In that event a quadratic
penalty suffices.
\end{proposition}
\begin{proof}
Taking $x\in N$ proves $A\succeq0$. For $z\in\ker A$ and fixed $w\in W$,
the value at $tz+w$ contains $2t\ip{z}{Dw}$, while $C(tz+w)=Cw$ is fixed.
Boundedness below for all real $t$ forces this coefficient to vanish.

Conversely, $\ker A\subseteq\ker D^*$ implies $D(W)\subseteq\operatorname{ran}A$.
Let $A^\dagger$ be the inverse on $\operatorname{ran}A$ and zero on $\ker A$.
Completing the square gives
\[
 \ip{n+w}{B(n+w)}=
 \ip{n+A^\dagger Dw}{A(n+A^\dagger Dw)}
 +\ip{w}{(G-D^*A^\dagger D)w}.
\]
The restriction $C|_W$ is injective; extend its inverse on $C(W)$ to a
linear map $T:F\to W$. The quadratic function
$\Phi(y)=\ip{Ty}{(D^*A^\dagger D-G)Ty}$ cancels the second term and vanishes
at zero. If $F$ has a chosen Euclidean norm, a sufficiently large nonnegative
multiple of $\|y\|^2$ also suffices, by finite-dimensional norm bounds.
\end{proof}
For a concrete counterexample let $H=\operatorname{diag}(0,2)$, $c=1$ and
$C(x,y)=x+y$. On $\ker C$, $\ip{x}{(H-I)x}=0$, and $C$ detects the negative
eigendirection of $H-I$. Nevertheless, on $(t,1-t)$ the repaired form is
$1-2t+\Phi(1)$, unbounded below for every finite $\Phi(1)$.
Kernel nonnegativity and low-mode detection alone do not suffice; cross-coupling
must be checked.

\section{Operator scaling and the actual consumer}
In the repository, the flat fixed-volume covariance endpoint uses
$K_0+aQ^*Q$ with the physical counting $\ell^2$ norm and the same unscaled
flat Hodge operator. Its coercivity constant is $\min(1,a)/C_P$, with $C_P$
from the full-periodic block-Poincare estimate. At fixed $d,N_c,N'$ the block
codomain is fixed and finite. Inspection therefore identifies the unscaled
flat base route with the setting of Theorem~\ref{thm:uniform}, with
$\Phi_L(y)=a_L\|y\|^2$ and no fine repair. Even allowing arbitrary scalar
$a_L$ or rescaling $Q_L$ does not yield a positive uniform full-space constant.
This is a source-level specialization in this text; no new Lean adapter to
that consumer is claimed.

The inspected endpoint is explicitly fixed-volume. Its positive rate may
depend on $L,N'$ and $C_P$, so there is no contradiction. Nor does the
obstruction exclude all possible uniform Combes--Thomas arguments. Its
locality radius and matrix-entry bounds also vary with $L$; coercivity alone
would not settle the rate. An interacting Hessian, a growing measurement
budget, a different sector norm or a physically rescaled operator must each
be audited separately.

In particular, writing $\ell=L+1$ so that $n_L=\ell N'$, replacing $K_0$
by $\ell^2K_0$ changes the band threshold exactly to
$4\pi^2dR^2/(N')^2$. The proof above
no longer drives it below every fixed $c>0$. This observation is an escape
from this proof, not an established coercivity theorem for the programme's
physical operator. No ultraviolet activity estimate, continuum limit,
Wilson-Hessian identification or Yang--Mills mass gap is obtained.

\section{Formal evidence and publication context}
The source-tree anchor containing all formal results used above is
\begin{center}\small\texttt{e31b428917aa767b4a1aad097db5cf92fb9cd23e}.\end{center}
The geometric and Fourier modules are ordinary proof dependencies, not
assumptions about a newly supplied band. Each development stage has complete
same-source fresh-clone Colab evidence. For the uniform endpoints the tested
source is
\begin{center}\small\texttt{17ca395d4e9de992b1c575454a112da9e6f97e86}.\end{center}
Its two runs
agree in compiled object, raw oracle log and source-manifest hashes. The
standalone budget lemma uses \texttt{propext, Quot.sound}; the other four
uniform-module reports also use \texttt{Classical.choice}. No extra axiom is
reported. Actual-log acceptance and three real-field rejection tests pass
under normal and optimized Python. These tests concern evidence handling,
not mathematical independence.

Short labels and an external declaration map replace long wrapped Lean names
in this manuscript. The companion source map records exact names, source SHAs,
and evidence directories. Reproduction commands remain in the existing
copyable notebooks rather than being split into misleading line fragments.
The fixed-map Schur criterion, its counterexample calculation and consumer
specialization above are analytic or source-inspection arguments; the whole
manuscript is not described as machine checked.

\clearpage
\subsection*{Model assessments and review context}
This text was prepared with the Codex agent in the author's ongoing research
chat. The exact served model identifier is not exposed in this preparation
record; a different named reviewer or reasoning setting is not inferred.
The agent had the prior manuscript objections, source audits and reproduction
records, and was not blind to the programme. No exact-PDF external model
assessment, numerical score, independent human peer review or editorial
acceptance is claimed for this text. Fresh runs of the same workflow are
reproducibility checks, not independent scientific assessments. Previous
version scores, if any, do not transfer to a revised artifact.

Lluis Eriksson is both the author and operator of AIRR.science. Any later
founder-author editorial decision must be recorded separately from model
assessments and from independent review. The existing record \cite{eriksson}
is an author-authorized historical import, not assessed under the current
AIRR screening. This preparation does not change its public version or
publication status. Code and manuscript licensing must be recorded separately;
no new license grant for the historical PDF is inferred.

\begin{thebibliography}{9}
\raggedright
\bibitem{eriksson} L. Eriksson, \emph{Exact-Kernel Fourier Families Obstruct
Every Penalty Factored Through a Block Map in a Flat Lattice Gauge Form},
ai.viXra 2608.0015v1 (2026), AIRR ARR-2026-7EPM3V3DP796BBRR v1.
\url{https://airr.science/papers/ARR-2026-7EPM3V3DP796BBRR/versions/v1/}.
Historical import, not peer reviewed or assessed by AIRR; checked 30 October 2026.
\bibitem{dodziuk} J. Dodziuk, \emph{Finite-difference approach to the Hodge
theory of harmonic forms}, American Journal of Mathematics \textbf{98}
(1976), 79--104. \url{https://doi.org/10.2307/2373615}.
Institutional source: \url{https://academicworks.cuny.edu/gc_pubs/180/}.
\end{thebibliography}
\end{document}
